\section{Background and related work}\label{sec:related}

\subsection{The Circum-Baltic contact area}\label{sec:related-cb}
Linguists have studied the languages around the Baltic as an area since \citet{dahl2001} introduced the term ``Circum-Baltic''. The area is unusual in that it is \emph{not} a Sprachbund definable by a shared feature bundle. \citet{koptjevskaja2001} conclude that it is better described as a \emph{contact superposition zone}: many intensive bilateral contacts among small groups of languages, layered over a long period, with no political unification under a single language. \citet{serzant2025} lists twenty-seven isoglosses running through the area; among the very few attested across all of it are a strong tendency toward initial stress (Scandinavian, Latvian, Žemaitian Lithuanian, all Finnic, German, some northern Russian dialects) and German lexical borrowings, ``found in all languages of the area with no exception.'' German during the Hanseatic period was ``probably the only lingua franca to have reached all parts'' of the area. Baltmeri's design reflects both facts: initial stress is adopted by default, and words already shared across families---overwhelmingly Hanseatic in origin---are inherited before any vote takes place (\S\ref{sec:procedure}). The general finding of contact linguistics that intensity of contact, not typological distance, predicts borrowing \citep{thomason1988} explains why a herring can be the same word in Danish and Polish while the two grammars remain alien to each other.

\subsection{The Hanseatic lexical legacy}\label{sec:related-hansa}
The scale of Middle Low German influence on the region is measurable. In Gellerstam's frequency study of the 6,000 most common Swedish words, 24\% are direct German loans and 30\% when German mediation is counted \citep{gellerstam1973,elmevik2012}; \citet{jahr1999} describes Hanseatic merchants settling in Bergen, Stockholm and Copenhagen and Low German entering everyday urban life. German loans are estimated at 22--25\% of Estonian stems \citep{hinderling1981}, and a World Loanword Database template applied to Finnish finds 26\% of the lexicon borrowed, 71\% of it Germanic and chiefly Swedish \citep{cronhamn2018}. Visby was ``the main centre of the Hanseatic League in the Baltic from the 12th to the 14th century'' \citep{unesco1995,dollinger1970}, which is why the Manual anchors its tie-break geography there: the sea's historical hub is the one point no modern state can claim.

\subsection{Zonal and naturalistic constructed languages}\label{sec:related-conlang}
Baltmeri belongs to the family of \emph{zonal} constructed languages, designed for a group of related or neighbouring languages rather than for the world \citep{gobbo2020}. Its closest methodological ancestor is IALA's Interlingua \citep{iala1951}, whose dictionary admitted a word ``when its presence is attested---in corresponding forms and with corresponding meanings---in at least three of the language units'' (Italian, French, English, Spanish/Portuguese, with German and Russian as substitutes), and rendered it in a \emph{prototype} form, ``the nearest common documented or hypothetical ancestor from which all contributing variants can be developed.'' Baltmeri's Hanse Rule is a threshold-plus-prototype rule of the same kind, with two differences: the threshold is counted over families rather than languages, and the form is not a reconstructed prototype but an attested stem chosen by vote. Interslavic is the second ancestor: its six Slavic sub-branches ``are weighed equally in establishing the largest common denominator \ldots\ to prevent any input language from getting undue weight'' \citep{merunka2019,vansteenbergen2016}, exactly the reasoning behind Baltmeri's one-family-one-vote rule; Interslavic's intelligibility surveys (mean 84\% cloze comprehension across 1,822 respondents) show that such a language can be passively understood by speakers who never learned it. Latin-derived and pan-Germanic projects \citep{libert2004} and the broad histories of \citet{okrent2009} and \citet{peterson2015} supply the wider context, from \citet{zamenhof1887} onward. What none of these projects attempted is full determinism: every one of them left the final form of a word to editorial judgement. Baltmeri's contribution to interlinguistics is to remove that judgement and to measure what happens when one does (\S\ref{sec:eval}).

\subsection{Computational comparison of vocabularies}\label{sec:related-comp}
The consonant skeleton at the heart of Baltmeri's procedure has a lineage in computational historical linguistics. \citet{dolgopolsky1986} proposed ten consonant classes such that sound changes within a class are more frequent than changes across classes; the Automated Similarity Judgment Program reduced IPA to 41 sound-class symbols and compared 40-item wordlists across thousands of languages \citep{brown2008,holman2008,wichmann2010,jaeger2018}; \citet{list2012,list2014} refined the classes for phonetic alignment and cognate detection, with accuracies approaching expert judgement \citep{list2017,rama2018}. Baltmeri's eleven classes (Table~\ref{tab:classes}) are coarser than any of these, deliberately: the aim is not to detect cognates but to detect \emph{shared words} regardless of whether they are cognate, borrowed or coincidental, since the sea does not care which. The basic-vocabulary tradition from \citet{swadesh1952,swadesh1955} to the Leipzig--Jakarta list \citep{tadmor2010,haspelmath2009} tells us where to expect agreement: roughly a quarter of an average lexicon is borrowed, but core meanings resist borrowing, so family votes on core vocabulary are least noisy while trade and sea vocabulary is where Hanse words cluster.

\subsection{The sea as subject; language and identity}\label{sec:related-sea}
\S\ref{sec:intro} has set out the ecological and legal context. Here we note the intellectual traditions the design draws on. Thalassology, from \citet{braudel1949} through \citet{horden2000,horden2006} to \citet{north2015} and \citet{kirby2000}, treats a sea as a historical subject with its own rhythms and connections rather than as a container for the nations around it; the medieval Baltic has been studied as an ``imagined community'' in exactly Anderson's sense \citep{jezierski2016,anderson2006}. The blue humanities \citep{mentz2009,gillis2012,steinberg2015} and posthuman feminist work on water \citep{neimanis2017,neimanis2017gotland} argue that the sea is not outside us. Rights-of-nature scholarship, from \citet{stone1972} to \citet{kauffman2021}, shows that legal personhood for ecosystems succeeds where it is embedded in coalitions and has human guardians who speak for the entity, as Te Pou Tupua speaks for the Whanganui \citep{teawatupua2017}. Ecolinguistics \citep{haugen1972,fill2001,stibbe2021} supplies the claim that the stories a language makes available shape ecological conduct; \citet{fishman1989} the link between language and self-definition. Baltmeri operationalises these positions: it gives the guardian something to speak.
